\section{Environment implementation}
\label{sec:implementation}

\subsection{The environment}
\label{sec:environment}

\slowlab{} is a harness: the components below are defined over a domain rather than within
one, and \S\ref{sec:domain} supplies the first domain that fills them in;
Figure~\ref{fig:overview} on page~\pageref{fig:overview} is the loop this
subsection specifies.


\subsubsection{Interaction protocol}
\label{sec:episode}

Figure~\ref{fig:protocol} is the whole interface. \texttt{validate\_design} is free
and unlimited; \texttt{submit\_design} commits the named slots for one cycle;
\texttt{advance\_to(day)} moves the simulation clock without resampling the crop, and
\texttt{observe} returns only measurements available on that day. Terminal outcomes are
released when the cycle completes. An interim recommendation records what the agent would
advise from the history currently visible. Budget is accounted in slot-cycles. The
recommendation is elicited, never consumed: it costs no budget, returns nothing, and
enters no later design --- an agent that revises it every round is not penalised for
having been wrong earlier. Design and recommendation are separate submissions, which is
what allows an agent to run a good campaign and then answer badly, and \S\ref{sec:eig} to
score the two apart.

\begin{figure}[t]
\centering
\footnotesize
\setlength{\fboxsep}{5pt}
%
\begin{minipage}[t]{0.225\linewidth}
\fbox{\begin{minipage}[t][45mm][t]{\dimexpr\linewidth-2\fboxsep-2\fboxrule}
\centering\textbf{\slowlab{} API}\par\vspace{3pt}
\raggedright\ttfamily\scriptsize
\pykw{class} \pycls{SlowLabEnv}:\\
\hspace*{1em}\pyfn{validate\_design}\\
\hspace*{1em}\pyfn{submit\_design}\\
\hspace*{1em}\pyfn{advance\_to}\\
\hspace*{1em}\pyfn{observe}\\
\hspace*{1em}\pyfn{observations}\\
\hspace*{1em}\pyfn{submit\_recommendation}\\[4pt]
\pykw{class} \pycls{Design}:\\
\hspace*{1em}treatments\\
\hspace*{1em}allocation\\
\hspace*{1em}randomization\_seed\\[4pt]
\pykw{class} \pycls{Task}:\\
\hspace*{1em}factors, lev(\,$\cdot$\,)\\
\hspace*{1em}R, |U\_r|, $\Delta$
\end{minipage}}
\end{minipage}
\hfill
\begin{minipage}[t]{0.335\linewidth}
\fcolorbox{black}{blue!4}{\begin{minipage}[t][45mm][t]{\dimexpr\linewidth-2\fboxsep-2\fboxrule}
\centering\textbf{One campaign}\par\vspace{3pt}
\raggedright\ttfamily\scriptsize
env = \pycls{SlowLabEnv}(task, seed)\\[2pt]
\pykw{for} r \pykw{in} \pykw{range}(R):\\
\hspace*{1em}D = agent.\pyfn{design}(env)\\
\hspace*{1em}env.\pyfn{validate\_design}(D)\\
\hspace*{1em}env.\pyfn{submit\_design}(D)\\
\hspace*{1em}env.\pyfn{advance\_to}(day)\\
\hspace*{1em}env.\pyfn{observe}(units, modality)\\
\hspace*{1em}x\_t = agent.\pyfn{recommend}(env)\\
\hspace*{1em}env.\pyfn{advance\_to}($\Delta$)\\[2pt]
env.\pyfn{submit\_recommendation}(x\_R)
\end{minipage}}
\end{minipage}
\hfill
\begin{minipage}[t]{0.385\linewidth}
\fbox{\begin{minipage}[t][45mm][t]{\dimexpr\linewidth-2\fboxsep-2\fboxrule}
\centering\textbf{What the environment does}\par\vspace{3pt}
\raggedright\scriptsize
the seed draws a \emph{site} $\theta\sim P_\Theta$, not a noise realisation\\[2.5pt]
validation is free; a rejection costs no campaign budget\\[2.5pt]
\textbf{the irreversible step}: $|\mathcal{U}_r|$ slots held for a full cycle\\[2.5pt]
advancing the clock exposes no future state or outcome\\[2.5pt]
observation returns a cached, time-appropriate measurement\\[2.5pt]
the visible history can update the recommendation\\[2.5pt]
episode end returns $\bar{\mathcal{R}}$ and $\mathcal{R}_{\mathrm{cum}}$
\end{minipage}}
\end{minipage}
\caption{\textbf{The Version~2 interaction protocol.} \emph{(left)} The design,
clock, measurement and recommendation calls available to an agent. \emph{(centre)} A
campaign may contain several inspection times within one committed cycle. \emph{(right)}
Committing occupies the units and cannot be undone; advancing exposes no future state;
observing reads the cached measurement for the requested day rather than drawing another
independent sample. Terminal outcomes remain hidden until $\Delta$.}
\label{fig:protocol}
\end{figure}

\subsubsection{Action space}
\label{sec:action}

A submission is a \texttt{Design}: named treatments, each a full vector of factor values; an
allocation mapping treatments to specific units; and a declared randomisation seed. The
agent chooses \emph{what} to test, \emph{how many units} each treatment receives, and
\emph{which} units those are---replication and allocation are decisions rather than
defaults. It may choose an observation day, modality and subset of units from the measurement
channels implemented by the task. Version~2 initially exposes cached sensor measurements
without a separate measurement charge. It cannot change treatment, terminate, clear or
replant a unit mid-cycle; a submitted design runs to completion.

\subsubsection{Facility and control levels}
\label{sec:facility}

A facility is $C$ blocks of $L$ units. Every factor carries a \emph{control level}: a
\textsc{block}-level factor takes one value for a whole block, a \textsc{unit}-level factor
is set per unit. The level is a physical property of the apparatus, not a modelling choice,
and a design assigning two different block-level values inside one block describes an
experiment that cannot be performed.

Two consequences follow, and we learned both by implementation. With $k$ block-level
factors, the number of distinct block-level \emph{combinations}---not levels per
factor---bounds the number of whole-plot treatments by $C$. And if a combination occupies
exactly one block, its effect is confounded with that block's, and unit-level replicates
cannot be spread. A valid design is therefore a \textbf{split-plot}: block-level factors as
whole plots replicated across at least two blocks, unit-level factors as subplots. This
structure is what makes allocation a decision, and it is absent from any environment in
which an experiment is a function call.

\subsubsection{Forward model and instances}
\label{sec:world}

The environment takes a published forward model $f_\theta$ mapping a factor vector to a mean
response, and a parameter distribution $P_\Theta$ describing how that model's coefficients
vary across sites in the literature. A seed selects an \emph{instance} by drawing
$\theta \sim P_\Theta$: a different problem, not the same problem observed again. The
response surface and the location of its optimum are therefore fixed by the published model
rather than by us.

Within an instance, each unit receives its own draw of the model's parameters rather than an
additive noise term. Unit-to-unit variation \emph{is} parameter variation, and modelling it
this way yields heteroscedastic, temporally correlated variability without further
assumptions. The nuisance effects of \S\ref{sec:conditions} are what an allocation must avoid
confounding with.

\subsubsection{Objective and cost}
\label{sec:costmodel}

The objective is a rate: revenue per unit area per unit time, less the full cost of the
resources the treatment consumed, $f_\theta(x) = \big(\text{revenue}(x) -
\text{cost}(x)\big) / \Delta$. Pricing is not decoration. Under a pure-yield objective a
forward model is typically monotone in its resource inputs, so their optima sit on the
boundary and those dimensions are trivial; charging for them is what produces interior
optima. Expressing the cost in the objective's own units is also what makes the
\emph{Costly} measure interpretable: a campaign's expenditure is directly comparable to
what it was trying to maximise.

The costs enter at weight one, and this is enforced rather than chosen. An earlier version
carried a multiplier $\lambda$ on the cost term, declared as a task-design parameter. It
was not one: it set where every optimum sat, and at $\lambda = 0.25$ it pushed four of six
factors onto their rails. The environment now raises if $\lambda \neq 1$
(Appendix~\ref{app:defects}). What \textsc{Transfer} moves mid-episode is a price in the
cost model, not a weight on it, so the objective stays the profit a grower would actually
bank in both halves of that task.

\subsubsection{Feasibility checking}
\label{sec:checking}

\texttt{validate\_design} checks \emph{mechanical} feasibility only---control levels, unit
conflicts, budget, factor ranges---and is free and unlimited; rejections cost no budget but
are counted. The design measure of \S\ref{sec:eig} is computed at scoring time and
cannot be called by the agent. Separating the two has a purpose: because feasibility is
free to check, a design that passes it yet is scientifically poor is unambiguously a
reasoning error rather than an interface error. Anything requiring a statistical
assumption---assumed variances, power targets---belongs to the tool layer, an experimental
condition to be ablated (Appendix~\ref{app:refs}), not a free service of the environment.

\subsubsection{Language-model interface}
\label{sec:harness}

A language model reaches the same interface through a harness that serialises the round. It
receives the facility layout and unit naming, the budget and rounds remaining, each factor's
range and control level, the hierarchy constraint, every time-stamped observation so far with the settings
that produced it, the current simulated day, the currently available units, and a JSON
schema. During an active cycle it may request a later observation day and measurement
channel or finish the cycle; it can update its recommendation from the resulting visible
history without changing the committed treatment.

The prompt states constraints and never gives advice: the words replication, randomisation,
blocking, split-plot, pure error and confounding do not appear, since those name the
capability under test, and a regression test enforces their absence. Failures are counted
separately---a reply that will not parse from a design that is mechanically infeasible---both
fed back verbatim and retried up to three times, and neither consuming budget, consistent with
\texttt{validate\_design} being free. ``Cannot emit valid JSON'' and ``cannot design an
experiment'' are therefore different numbers. Exhausting the retries forfeits that round and
the campaign continues, so one unparseable reply is not as costly as a failed campaign.

\subsubsection{Configuration and extension}
\label{sec:interfaces}

Configurations are declarative: factors and their control levels, facility shape, rounds,
per-round width, variance components, cost weight. Adding one requires no environment code.
Agents implement a single \texttt{run(env)} method and are registered by name; the reference
strategies of Appendix~\ref{app:refs} and the harness above use the same interface, and the
harness needs only a callable mapping a message list to a string, so no vendor SDK is
required. A new domain supplies $f_\theta$, $P_\Theta$, a cost function in the
objective's units, and a facility whose units are not interchangeable
(Appendix~\ref{app:admit}).
